% 349_SU2L_Chiralitaet_Galois_De_ch.tex
% Johann Pascher, 4. September 2026

\providecommand{\BM}{\textbf{[B]}}
\providecommand{\KM}{\textbf{[K]}}
\providecommand{\SM}{\textbf{[S]}}

\phantomsection
\addcontentsline{toc}{chapter}{SU$(2)_L$-Chiralität aus der Galois-Struktur}

\setcounter{section}{0}
\section*{Statuszeichen}
\begin{center}
\begin{tabular}{@{}lp{10.5cm}@{}}
\textbf{[K]} & \emph{Konfirmiert} --- algebraisch vollständig. \\[4pt]
\textbf{[B]} & \emph{Belegt} --- durch Prüfskript gesichert. \\[4pt]
\textbf{[S]} & \emph{Spekulativ} --- Vermutung, noch offen. \\
\end{tabular}
\end{center}

\section*{Abstract}

Die SU$(2)_L$-Chiralitätsstruktur des Standardmodells folgt algebraisch
aus der Cl$(6)$-Gradparität, die bereits in Dok.~341 (Z$_3$-Trennung der
geraden und ungeraden Vakua) etabliert ist. Der Pseudoskalar
$\gamma_5 = e_1 e_2 e_3 e_4 e_5 e_6 \in \mathrm{Cl}(6)$ hat Eigenwert
$(-1)^r$ auf Cl$(6)$-Grad~$r$. Die Su/Sd-Trennung (jE7-Projektor, Dok.~344)
ist deshalb exakt die Rechts/Links-Aufspaltung: Su = gerade Grade =
rechtshändig \BM, Sd = ungerade Grade = linkshändig \BM.
Die offene Frage aus Dok.~347/348 ist damit geschlossen \BM.

\section{Ausgangslage}

Dok.~341 zeigt (Abschn.\ Z$_3$-Trennung, \BM): die Witt-Paar-Rotation
$\sigma$ trennt die Vakua in zwei Klassen:
\begin{itemize}
\item \textbf{Gerade Vakua} $\mathrm{Vss}[0,2,4]$: einzelne
  $\mathbb{Z}_3$-Fixpunkte; entsprechen geraden Cl$(6)$-Graden.
\item \textbf{Ungerade Vakua} $\mathrm{Vss}[1,3,5]$: ein
  $\mathbb{Z}_3$-Dreier-Orbit; entsprechen ungeraden Cl$(6)$-Graden.
\end{itemize}
Dok.~344 zeigt (\BM): Su = Furey-Fermionen des Up-Typs (N Witt-Paare besetzt,
$N$ gerade); Sd = Down-Typ ($N$ ungerade). Der jE7-Projektor trennt beide.

\section{Satz A --- Gradparität und $\gamma_5$-Eigenwert}

\begin{theorem}[$\gamma_5$-Eigenwert nach Cl$(6)$-Grad \BM]
Im Clifford-Raum $\mathrm{Cl}(6)$ gilt für jedes Grad-$r$-Element $X$:
\[
  \gamma_5 \cdot X = (-1)^r \cdot X, \qquad
  \gamma_5 = e_1 e_2 e_3 e_4 e_5 e_6 \in \mathrm{Grad}\,6.
\]
\textbf{Gerade Grade} ($r = 0,2,4,6$): $\gamma_5 X = +X$ \emph{(rechtshändig)}.\\
\textbf{Ungerade Grade} ($r = 1,3,5$): $\gamma_5 X = -X$ \emph{(linkshändig)}.
\end{theorem}

\begin{proof}
$\gamma_5$ antikommutiert mit jedem Generator $e_i$ (Grad~1),
da $\{e_i, e_j\} = 2\delta_{ij}$ und $\gamma_5$ enthält alle sechs Generatoren.
Für Grad-$r$-Blade $X = e_{i_1}\cdots e_{i_r}$:
$\gamma_5 X = (-1)^r X\,\gamma_5$, also $\gamma_5 X = (-1)^r X$
als $\gamma_5^2 = +1$.
\end{proof}

\section{Satz B --- Su = gerade Grade = rechtshändig}

\begin{theorem}[Su-Sektor ist rechtshändig \BM]
Die Su-Zustände (Dok.~344) liegen in geraden Cl$(6)$-Graden:
\begin{center}
\begin{tabular}{llll}
\toprule
Zustand & $N$ & Cl$(6)$-Grad & Chiralität \\
\midrule
$\nu$ (Su[0]) & 0 & 0 (Skalar) & rechtshändig \\
$\bar u_d$ (Su[1]) & 1 & 2 (Bivektor) & rechtshändig \\
$u$ (Su[2]) & 2 & 4 & rechtshändig \\
$e_R$ (Su[3]) & 3 & 6 (Pseudosk.) & rechtshändig \\
\bottomrule
\end{tabular}
\end{center}
$\gamma_5$-Eigenwert $= +1$ für alle Su-Zustände.
\end{theorem}

\section{Satz C --- Sd = ungerade Grade = linkshändig}

\begin{theorem}[Sd-Sektor ist linkshändig \BM]
Die Sd-Zustände liegen in ungeraden Cl$(6)$-Graden:
\begin{center}
\begin{tabular}{llll}
\toprule
Zustand & $N$ & Cl$(6)$-Grad & Chiralität \\
\midrule
$\bar\nu$ (Sd[0]) & 0 & 1 & linkshändig \\
$d$ (Sd[1]) & 1 & 3 & linkshändig \\
$\bar d$ (Sd[2]) & 2 & 5 & linkshändig \\
$e_L$ (Sd[3]) & 3 & 7* & linkshändig \\
\bottomrule
\end{tabular}
\end{center}
$\gamma_5$-Eigenwert $= -1$ für alle Sd-Zustände. (*Grad 7 mod 8 in Cl(6) als Multigrad.)
\end{theorem}

\section{Satz D --- SU$(2)_L$-Dubletts aus dem Sd-Sektor}

\begin{theorem}[SU$(2)_L$-Dublett-Struktur aus Gradparität \BM]
Die linkshändigen Dubletts des Standardmodells sind genau die
Sd-Zustände (ungerade Cl$(6)$-Grade); die rechtshändigen Singuletts
sind genau die Su-Zustände (gerade Grade):
\[
  \text{Dublett: } (\nu_L, e_L) = (\text{Sd}[0], \text{Sd}[3]), \quad
  (\bar u_L, d_L) = (\text{Sd-Quarks}),
\]
\[
  \text{Singlett: } e_R = \text{Su}[3],\quad
  u_R = \text{Su}[2],\quad
  \nu_R = \text{Su}[0] \text{ (Dirac-Fall)}.
\]
SU$(2)_L$ wirkt auf den Sd-Sektor (linkshändig) und
nicht auf den Su-Sektor (rechtshändig).
\end{theorem}

\textbf{Verbindung zu Dok.~341 \BM.} Die Z$_3$-Fixpunkte (gerade Vakua
$[0,2,4]$) sind identisch mit dem Su-Sektor (rechtshändig). Der
Z$_3$-Dreier-Orbit (ungerade Vakua $[1,3,5]$) ist identisch mit dem
Sd-Sektor (linkshändig). Die Chiralitätstrennung ist damit algebraisch
erzwungen durch die $\mathbb{Z}_3$-Symmetrie des Torus, nicht postuliert.

\section{Satz E --- Z$_3$-Fixpunkt/Orbit = Chiralität}

\begin{theorem}[$\mathbb{Z}_3$-Struktur erzwingt Chiralitätstrennung \BM]
Die $\sigma$-Rotation aus Dok.~341 erzeugt:
\[
  \text{Vss}[0,2,4] \;(\sigma\text{-Fixpunkte}) = \text{rechtshändig},
\]
\[
  \text{Vss}[1,3,5] \;(\sigma\text{-Dreier-Orbit}) = \text{linkshändig}.
\]
Die SU$(2)_L$-Chiralitätstrennung ist eine direkte Konsequenz der
$\mathbb{Z}_3$-Symmetrie des T$^4$/Z$_3$-Torus.
\end{theorem}

\section{Zusammenfassung}

\begin{table}[H]\centering
\caption{Ergebnisse Dok.~349}
\begin{tabular}{lll}\toprule
Satz & Inhalt & Status \\\midrule
A & $\gamma_5$-Eigenwert $(-1)^r$ auf Cl$(6)$-Grad $r$ & \BM \\
B & Su = gerade Grade = rechtshändig & \BM \\
C & Sd = ungerade Grade = linkshändig & \BM \\
D & SU$(2)_L$-Dubletts = Sd-Sektor & \BM \\
E & Z$_3$-Fixpunkte/Orbit = rechts/links & \BM \\\bottomrule
\end{tabular}
\end{table}

\textbf{Was geschlossen ist.} Die in Dok.~347/348 als offen gekennzeichnete
SU$(2)_L$-Chiralitätsstruktur ist damit aus der Galois/Clifford-Algebra
hergeleitet \BM. Kein neues Postulat: die Gradparität war in Dok.~341 bereits
etabliert; der Schritt ist deren explizite physikalische Auswertung.

\textbf{Was noch offen bleibt \SM.} Die exakten CKM-Mischungswinkel und die
CP-Phase $\delta_{CP}$ (Dok.~348). Die Brücke liegt jetzt vollständiger:
GF$(27)^*$ (Galois-Orbits) → Cl$(6)$ (Gradstruktur) → $\gamma_5$ (Chiralität)
→ SU$(2)_L$ (Dublett/Singlett) ist jetzt \BM.

\section*{Prüfskript}
\texttt{pruef\_349\_chiralitaet.py} --- Sätze A--E, alle Assertions bestanden \BM.

\section*{Vermerk zur Verwendung von KI-Assistenz}
Erstellt mit Unterstützung von Claude (Anthropic).
Physikalische Interpretationen und Schlussfolgerungen:
Johann Pascher (ORCID 0009-0000-6518-4064).

\begin{thebibliography}{9}
\bibitem{P341} J.~Pascher, \textit{Dok.~341: GF(27) in GALG und FFGFT}, FFGFT-Korpus (Sept.\ 2026).
\bibitem{P344} J.~Pascher, \textit{Dok.~344: Furey-Fermionen in GALG und FFGFT}, FFGFT-Korpus (Sept.\ 2026).
\bibitem{P347} J.~Pascher, \textit{Dok.~347: Gell-Mann-Nishijima aus Galois}, FFGFT-Korpus (Sept.\ 2026).
\bibitem{P348} J.~Pascher, \textit{Dok.~348: Generationenstruktur aus GF$(27)$}, FFGFT-Korpus (Sept.\ 2026).
\bibitem{Furey2015} C.~Furey, Ph.D.\ Thesis, Waterloo 2015, arXiv:1611.09182.
\end{thebibliography}
